
\section{一个「\trUniversal」\trVariationalPrinciple}
\label{sec:universal-variational-principle}
\hzsync{theorem}{4}%
\hzsync{lemma}{2}%
\hzsync{equation}{17}%

5.1 节中的结果基于一个\trVariationalPrinciple，我们现在%
便转向这个原理。我们的目标是建立\trActionFunctional
\[
a_{H}(x)=a(x)-b_{H}(x),\qquad x\in E
\]
的一个特殊\trCriticalPoint{}的存在性；这个\trActionFunctional{}已在第 3 章%
中引入。与那里所用的原理不同，我们将设计一个不依赖于 $H\in\mathcal{H}$ 之%
选取的\trMinimax{}原理。回顾：$E$ 是\trHilbert{}空间 $E=\Hh{1}{2}\bigl(S^{1},\R^{2n}\bigr)$，%
它分裂成一个\trOrthogonalSum{} $E=E^{-}\oplus E^{0}\oplus E^{+}$；并且，若%
$x=x^{-}+x^{0}+x^{+}$，则
\[
a(x)=\frac{1}{2}\bigl\lVert x^{+}\bigr\rVert^{2}-\frac{1}{2}\bigl\lVert x^{-}\bigr\rVert^{2}
\]
以及
\[
b_{H}(x)=\int_{0}^{1}H\bigl(t,x(t)\bigr)\ dt.
\]
正如第 3 章中那样，人们可以验证：\trHamiltonianEquation{} $\dot{x}=X_{H}(x)$%
的特殊解 $x\in C^{\infty}\bigl([0,1],\R^{2n}\bigr)$
\[
x(t)=\varphi_{H}^{t}(x_{0})\quad\text{以及}\quad x_{0}\in\operatorname{Fix}(\varphi_{H})
\]
与 $a_{H}$ 在 $E$ 上的诸\trCriticalPoint{}相一致。我们将挑出由子集 $F\subset E$%
组成的一个族 $\mathcal{F}$，使得对每个 $H\in\mathcal{H}$
\[
\gamma(H):=\sup_{F\in\mathcal{F}}\ \inf_{x\in F}\ a_{H}(x)
\]
是一个实数，并且是 $a_{H}$ 的一个\trCriticalValue。关键之处将在于：$\mathcal{F}$%
不依赖于 $H\in\mathcal{H}$ 的选取，因而对 $\mathcal{H}$ 而言是一个「\trUniversal」%
\trMinimax{}集。这在后面出于如下理由将对我们颇为重要：每当手头有一个%
\trFunctional{}族连同「\trUniversal」\trMinimax{}集 $\mathcal{F}$，那么从这个族中%
两个\trFunctional{}之间的逐点\trEstimate{} $\Phi\leq\Psi$，便可对相应的%
\trMinimax{}值推出同样的不等式。注意
\[
a_{H}(x^{+})\to+\infty\quad\text{当}\ \lVert x^{+}\rVert\to\infty
\]
其中 $x^{+}\in E^{+}$；而
\[
a_{H}(x^{-}+x^{0})\leq M\quad\text{对一切}\ x^{-}+x^{0}\in E^{-}\oplus E^{0}
\]
其中常数 $M$ 依赖于 $H$。

为了定义 $\mathcal{F}$，我们先引入一个特定的 $E$ 之\trHomeomorphism{}群，它的%
提出是受 $a_{H}$ 的诸\trGradient\trFlow{}之结构启发的。

\begin{definition}\label{def:homeomorphism-group-G}
我们如下定义 $E$ 的\trHomeomorphism{}之集 $G$。一个\trHomeomorphism{} $h$%
属于 $G$，如果 $h$ 与 $h^{-1}$ 都把\trBounded{}集映成\trBounded{}集，并且存在%
连续映射 $\gamma^{\pm}:E\to\R$ 与 $k:E\to E$，它们把\trBounded{}集映入%
预\trCompact{}集，此外还存在 $r=r(h)$ 使得
\[
k(x)=0,\qquad\gamma^{\pm}(x)=0
\]
对所有满足 $\lVert x\rVert\geq r$ 的 $x\in E^{+}$ 成立。进而，
\[
h(x)=e^{\gamma^{+}(x)}x^{+}+x^{0}+e^{\gamma^{-}(x)}x^{-}+k(x)
\]
对一切 $x\in E$ 成立。
\end{definition}

\begin{lemma}\label{lem:g-is-group}
$G$ 是一个群。
\end{lemma}

\begin{proof}
我们必须证明：与 $h,g\in G$ 一起，也有 $h^{-1}$ 与 $h\circ g\in G$。由对 $h$%
的假设我们有
\begin{align*}
e^{-\gamma^{\pm}}P^{\pm}h(x)\ &=\ x^{\pm}+e^{-\gamma^{\pm}}P^{\pm}k(x)\\
P^{0}h(x)\ &=\ x^{0}+P^{0}k(x).
\end{align*}
于是
\begin{align*}
x\ &=\ x^{-}+x^{0}+x^{+}\\
&=\ e^{-\gamma^{-}(x)}P^{-}h(x)+P^{0}h(x)+e^{-\gamma^{+}(x)}P^{+}h(x)\\
&\quad-\bigl(e^{-\gamma^{-}(x)}P^{-}+P^{0}+e^{-\gamma^{+}(x)}P^{+}\bigr)k(x).
\end{align*}
记 $x=h^{-1}(y)$ 及 $y=h(x)$，则对 $y=y^{-}+y^{0}+y^{+}$ 得到
\begin{align*}
h^{-1}(y)\ &=\ e^{-\gamma^{-}(h^{-1}(y))}y^{-}+y^{0}+e^{-\gamma^{+}(h^{-1}(y))}y^{+}\\
&-\ \bigl(e^{-\gamma^{-}(h^{-1}(y))}P^{-}+P^{0}+e^{-\gamma^{+}(h^{-1}(y))}P^{+}\bigr)k\bigl(h^{-1}(y)\bigr)\\
&=:\ e^{\sigma^{-}(y)}y^{-}+y^{0}+e^{\sigma^{+}(y)}y^{+}+\hat{k}(y).
\end{align*}
利用 $\gamma^{\pm}$ 与 $k$ 的性质，再结合 $h^{-1}$ 把\trBounded{}集映入%
\trBounded{}集这一事实，我们得出 $\sigma^{\pm}$ 与 $\hat{k}$ 具有所要的性质，%
从而 $h^{-1}\in G$。类似地，可以验证 $h\circ g\in G$。
\end{proof}

这个由 $E$ 的诸子集组成的特定族 $\mathcal{F}$ 定义为
\[
\mathcal{F}=\bigl\{\,h(E^{+})\ \bigm|\ h\in G\,\bigr\}.
\]
$G$ 的关键\trTopology{}性质蕴含在下述直观明了的相交结果之中：

\begin{lemma}\label{lem:intersection-nonempty}
若 $h\in G$，则
\[
h(E^{+})\cap\bigl(E^{-}\oplus E^{0}\bigr)\neq\varnothing.
\]
\end{lemma}

\begin{proof}
我们需要求得一个满足 $P^{+}h(x)=0$ 的 $x\in E^{+}$。这等价于%
$e^{\gamma^{+}(x)}x+P^{+}k(x)=0$ 且 $x\in E^{+}$，亦即
\[
x=-e^{-\gamma^{+}(x)}P^{+}k(x)=:T(x),\quad\text{且}\ x\in E^{+}\ .
\]
\trOperator{} $T:E^{+}\to E^{+}$ 把 $E^{+}$ 映入 $E^{+}$ 的一个预\trCompact{}%
子集，从而对充分大的\trBall{} $B$ 有 $T:\overline{B}\to\overline{B}$，于是%
\trSchauder{}不动点定理保证了一个解的存在，引理得证。
\end{proof}

\begin{theorem}\label{thm:minimax-critical-value}
设 $H\in\mathcal{H}$，则由
\[
\gamma(H):=\sup_{F\in\mathcal{F}}\ \inf\ a_{H}(F),
\]
所定义的 $\gamma(H)$ 是 $a_{H}$ 的一个\trCriticalValue。
\end{theorem}

\noindent\begin{minipage}{\linewidth}
\centering
\begin{tikzpicture}[scale=0.9]
  \draw[dotted] (4.5,6.55) -- (4.5,6.05);
  \draw (4.5,6.05) -- (4.5,1.7);
  \draw (4.5,1.7) -- (4.5,1.05);
  \draw[dotted] (4.5,1.05) -- (4.5,0.6);
  \draw (0.35,3.95) -- (2.75,4.6) -- (8.75,3.95) -- (6.35,3.3) -- cycle;
  \draw (4.5,5.9) .. controls (4.8,5.1) and (5.35,4.35) .. (5.4,3.6);
  \draw (5.4,3.6) .. controls (5.4,2.95) and (4.8,2.35) .. (4.5,1.7);
  \fill (5.4,3.6) circle (2pt);
  \draw (5.02,4.9) -- (5.75,5.25);%
  \node[anchor=south west] at (5.75,5.25) {$h(E^{+})$};
  \node[anchor=east] at (4.38,6.28) {$E^{+}$};
  \node[anchor=east] at (4.35,3.62) {$O$};
  \node at (6.7,3.88) {$E^{-}\oplus E^{0}$};
\end{tikzpicture}

图 5.1
\end{minipage}

本节其余部分将用于这个存在性定理的证明。我们将分几步进行，首先从相交引理%
推出 $\gamma(H)$ 是一个实数。定义
\begin{align*}
q^{+}(H)\ &=\ \int_{0}^{1}\sup_{x}H(t,x)\ dt\\
q^{-}(H)\ &=\ \int_{0}^{1}\inf_{x}H(t,x)\ dt.
\end{align*}
于是 $q^{+}(H)\geq0$ 且 $q^{-}(H)\leq0$，因为 $H\in\mathcal{H}$。

\begin{lemma}\label{lem:gamma-real}
对每个 $H\in\mathcal{H}$
\[
-q^{+}(H)\leq\gamma(H)\leq-q^{-}(H).
\]
\end{lemma}

\begin{proof}
设 $F\in\mathcal{F}$，则由\cref{lem:intersection-nonempty}有%
$F\cap\bigl(E^{-}\oplus E^{0}\bigr)\neq\varnothing$，于是我们可以估计
\begin{align*}
\inf a_{H}(F)\ &\leq\ \sup a_{H}\bigl(E^{-}\oplus E^{0}\bigr)\\
&\leq\ \sup_{x\in E}\bigl[-b_{H}(x)\bigr]\\
&\leq\ -\int_{0}^{1}\inf_{x}H(t,x)\ dt.
\end{align*}
由于
\[
\inf a_{H}(E^{+})\ \geq\ -\int_{0}^{1}\sup_{x}H(t,x)\ dt\ ,
\]
引理得证。
\end{proof}

若 $\gamma(H)=0$，则 $\gamma(H)$ 是一个\trCritical{}水平，它特别地含有 $H$ 的%
\trSupport{}之外的全部常值解。因此，我们只须证明 $\gamma(H)\neq0$ 是 $a_{H}$%
的一个\trCriticalValue。为此我们需要下述\trCompactness。

\begin{lemma}\label{lem:ps-gamma}
设 $\gamma(H)\neq0$，并设序列 $x_{j}\in E$ 满足
\[
\nabla a_{H}(x_{j})\to0,\qquad\text{且}\ a_{H}(x_{j})\to\gamma(H).
\]
则 $x_{j}$ 具有一个收敛的子序列。
\end{lemma}

\begin{proof}
我们回顾第 3 章中的
\[
\nabla a_{H}(x)=x^{+}-x^{-}-\nabla b_{H}(x),
\]
其中 $\nabla b_{H}:E\to E$ 是全局\trLipschitz{}连续的，并把\trBounded{}集映入%
预\trCompact{}集；此外 $\nabla b_{H}(x)=j^{*}\nabla H(x)$，其中 $\nabla H(x)$%
是 $b_{H}$ 相对于 $L^{2}$ 的\trGradient。

设序列 $x_{j}\in E$ 在 $E$ 中\trBounded。于是 $\nabla b_{H}(x_{j})$ 预\trCompact，%
取一个子序列后不妨设 $\nabla b_{H}(x_{j})\to y$（于 $E$ 中）。由假设%
$\nabla a_{H}(x_{j})\to0$，故 $x_{j}^{+}-x_{j}^{-}\to y$。此外，由于%
$\dim E^{0}<\infty$，我们可以找到一个子序列使 $x_{j}^{0}$ 也收敛，从而可见%
$x_{j}$ 具有收敛的子序列。

因此，我们只须证明序列 $x_{j}$ \trBounded。用反证法，取子序列后可设%
$\lVert x_{j}\rVert\to\infty$。由于 $H\in\mathcal{H}$，序列 $\nabla H(x_{j})$ 在%
$L^{2}$ 中\trBounded，故 $\nabla b_{H}(x_{j})=j^{*}\nabla H(x_{j})$ 在 $E$ 中预%
\trCompact。于是，取一个子序列，$x_{j}^{+}-x_{j}^{-}\to y$（于 $E$ 中，同时也在%
$L^{2}$ 中）。这意味着 $\lvert x_{j}^{0}\rvert\to\infty$。从而我们可以找到一个%
子序列使得逐点地
\[
\lvert x_{j}(t)\rvert\to\infty,\quad\text{对\trAlmostEvery{}的}\ t
\]
成立，这意味着对\trAlmostEvery{}的 $t$ 有 $H\bigl(t,x_{j}(t)\bigr)\to0$。于是由%
\trLebesgue{}定理，$b_{H}(x_{j})\to0$。又由于对\trAlmostEvery{}的 $t$ 还有%
$\nabla H\bigl(t,x_{j}(t)\bigr)\to0$，故 $\nabla H(x_{j})\to0$（于 $L^{2}$ 中），%
从而 $\nabla b_{H}(x_{j})\to0$（于 $E$ 中）。由 $\nabla a_{H}(x_{j})\to0$ 我们%
现在得出 $x_{j}^{+}-x_{j}^{-}\to0$。于是
\[
a_{H}(x_{j})=a(x_{j})-b_{H}(x_{j})\to0.
\]
这与 $a_{H}(x_{j})\to\gamma(H)\neq0$ 的假设相矛盾，并表明序列 $x_{j}$%
\trBounded。引理的证明至此完毕。
\end{proof}

我们现已可以结束\cref{thm:minimax-critical-value}的证明了，我们把它重新%
表述为

\begin{lemma}\label{lem:gamma-critical-value}
实数 $\gamma(H)$ 是 $a_{H}$ 的一个\trCriticalValue，这对每个 $H\in\mathcal{H}$%
都成立。
\end{lemma}

\begin{proof}
回顾：由\cref{lem:gamma-real}，$\gamma(H)$ 是一个实数。若 $\gamma(H)=0$，则%
无可再证，我们遂设 $\gamma(H)\neq0$。此时，\cref{lem:ps-gamma}中所表述的%
那个\trPalaisSmale{}条件成立。用反证法，我们设 $\gamma(H)$ 不是 $a_{H}$ 的%
\trCriticalValue。于是，依据\cref{lem:ps-gamma}，可以找到一个 $\varepsilon>0$%
满足
\begin{equation}\label{eq:gradient-estimate}
\bigl\lvert\nabla a_{H}(x)\bigr\rvert\ \geq\ \varepsilon,\quad%
\text{若}\ \ a_{H}(x)\in\bigl[\gamma(H)-\varepsilon,\ \gamma(H)+\varepsilon\bigr]
\end{equation}
为了先使\trGradient\trFlow{}减速，我们取一个\trSmooth{}映射%
$\hat{\sigma}:\R\to[0,1]$，满足
\begin{align*}
\hat{\sigma}(t)\ &=\ 1,\quad\text{若}\ t\in\bigl[\gamma(H)-\varepsilon,\gamma(H)+\varepsilon\bigr]\\
\hat{\sigma}(t)\ &=\ 0,\quad\text{若}\ t\in\R\setminus\bigl[\gamma(H)-2\varepsilon,\gamma(H)+2\varepsilon\bigr]
\end{align*}
并定义 $\sigma(x)=\hat{\sigma}\bigl(a_{H}(x)\bigr)$（$x\in E$）。修改后的%
\trGradient{}方程
\begin{equation}\label{eq:modified-gradient}
\dot{x}=\sigma(x)\nabla a_{H}(x)=:V(x)\quad\text{于}\ E\ \text{上}
\end{equation}
的\trFlow{}对一切时间都存在，我们将把它记作
\[
(s,x)\mapsto x\cdot s:\R\times E\to E\ .
\]
借助这个\trFlow，我们取 $T=\frac{2}{\varepsilon}$，并由
\begin{equation}\label{eq:flow-homeomorphism}
h(x)=x\cdot T,\qquad x\in E.
\end{equation}
定义 $E$ 的一个\trHomeomorphism{} $h$。由 $\gamma(H)$ 的定义，存在%
$F\in\mathcal{F}$ 满足
\[
F\subset\bigl\{\,x\ \bigm|\ a_{H}(x)\geq\gamma(H)-\varepsilon\,\bigr\},
\]
而我们将要证明
\begin{equation}\label{eq:hF-inclusion}
h(F)\subset\bigl\{\,x\ \bigm|\ a_{H}(x)\geq\gamma(H)+\varepsilon\,\bigr\}.
\end{equation}
把证明暂且搁下，我们先展示引理是如何收尾的。由 $\mathcal{F}$ 的定义我们有%
$F=g(E^{+})$。于是，若 $h\in G$（这个事实将在后面得到证明），则由于 $G$%
是一个群（\cref{lem:g-is-group}），亦有 $h\circ g\in G$。从而%
$h(F)=h\circ g(E^{+})\in\mathcal{F}$，于是我们得到矛盾
\[
\gamma(H)\ \geq\ \inf a_{H}\bigl(h(F)\bigr)\ \geq\ \gamma(H)+\varepsilon,
\]
引理随之得证。为了验证 \eqref{eq:hF-inclusion}，我们证明
\begin{equation}\label{eq:implication-shift}
a_{H}(x)\geq\gamma(H)-\varepsilon\ \Longrightarrow\ a_{H}\bigl(h(x)\bigr)\geq\gamma(H)+\varepsilon.
\end{equation}
用反证法，我们设 $a_{H}\bigl(h(x)\bigr)<\gamma(H)+\varepsilon$。由于%
$s\mapsto a_{H}(x\cdot s)$ 是单调递增的，我们有
\[
\gamma(H)+\varepsilon\ >\ a_{H}\bigl(h(x)\bigr)\ \geq\ a_{H}(x\cdot s)\ \geq\ a_{H}(x)\ \geq\ \gamma(H)-\varepsilon
\]
对一切 $0\leq s\leq T$ 成立。于是 $a_{H}(x\cdot s)\in[\gamma(H)-\varepsilon,\gamma(H)+\varepsilon]$，%
从而 $\sigma(x\cdot s)=1$（$0\leq s\leq T$）。我们看到：在这个时间区间内，%
$x(s):=x\cdot s$ 求解\trGradient{}方程 $\dot{x}=\nabla a_{H}(x)$。于是，对这个%
\trGradient{}使用\trEstimate{} \eqref{eq:gradient-estimate}：
\begin{align*}
\gamma(H)+\varepsilon\ &>\ a_{H}\bigl(h(x)\bigr)=a_{H}(x\cdot T)\\
&=\ a_{H}(x)+\int_{0}^{T}\bigl\lvert\nabla a_{H}(x\cdot s)\bigr\rvert^{2}\,ds\\
&\geq\ \gamma(H)-\varepsilon+T\varepsilon^{2}=\gamma(H)+\varepsilon,
\end{align*}
这是一个矛盾。\eqref{eq:implication-shift} 已获验证，剩下的还须证明%
$h\in G$。

回顾
\begin{equation}\label{eq:V-formula}
V(x)=-\sigma(x)x^{-}+\sigma(x)x^{+}-\sigma(x)\nabla b_{H}(x).
\end{equation}
为了给 $\dot{x}=V(x)$ 在 $E$ 上的\trFlow{}找到一个表示公式，我们对取定的%
$x\in E$ 考察方程
\begin{equation}\label{eq:linear-inhom}
\frac{d}{ds}u=A_{x}(s)u+f_{x}(s)\quad\text{于}\ E\ \text{上},
\end{equation}
其中 $A_{x}(s)=-\sigma(x\cdot s)P^{-}+\sigma(x\cdot s)P^{+}$ 是 $E$ 上一个%
\trTimeDependent{}但\trLinear{}的\trVectorField，而%
$f_{x}(s)=-\sigma(x\cdot s)\nabla b_{H}(x\cdot s)$。若 $u(s)$ 是%
\eqref{eq:linear-inhom} 具初值 $u(0)=x$ 的一个解，则由常%
\trDifferentialEquation{}的唯一性定理知 $u(s)=x\cdot s$。\eqref{eq:linear-inhom}%
之解的\trVariationOfConstantsFormula{}现在便给出方程 $\dot{x}=V(x)$ 之%
\trFlow{}的如下表示：\footnote{原书印作「$-P^{0}$」；因 $A_{x}(s)$ 在 $E^{0}$ 上为零，其演化算子在 $E^{0}$ 上应为恒等，且与上行表示中 $x^{0}$ 之系数相一致，中间项应为 $+P^{0}$，此处径改。——译者注}
\begin{align*}
x\cdot t\ &=\ e^{\int_{0}^{t}-\sigma(x\cdot s)ds}x^{-}+x^{0}+e^{\int_{0}^{t}\sigma(x\cdot s)ds}x^{+}\\
&\quad-\int_{0}^{t}\Bigl[e^{\int_{s}^{t}-\sigma(x\cdot\tau)d\tau}P^{-}+P^{0}+e^{\int_{s}^{t}\sigma(x\cdot\tau)d\tau}P^{+}\Bigr]\sigma(x\cdot s)\nabla b_{H}(x\cdot s)ds\\
&=:\ e^{\gamma^{-}(x)}x^{-}+x^{0}+e^{\gamma^{+}(x)}x^{+}+k(x).
\end{align*}
我们将证明：对每个 $t\in\R$，这个\trFlow{}映射都属于群 $G$。由于%
\trVectorField{} $V$ 是全局\trLipschitz{}连续的，它把 $E$ 中的\trBounded{}集%
映入\trBounded{}集；当然，这对逆映射同样成立。注意，存在 $r>0$，使得当%
$x\in E^{+}$ 满足 $\lVert x\rVert\geq r$ 时 $a_{H}(x)\geq\gamma(H)+2\varepsilon$，%
从而 $\sigma(x)=0$，进而 $V(x)=0$。于是，对一切 $s\in\R$ 与一切满足%
$\lVert x\rVert\geq r$ 的 $x\in E^{+}$ 都有 $\sigma(x\cdot s)=0$。我们得出：对%
满足 $\lVert x\rVert\geq r$ 的 $x\in E^{+}$ 有 $\gamma^{\pm}(x)=0$ 且 $k(x)=0$。%
最后，映射 $k$ 的所要的\trCompactness{}由映射 $\nabla b_{H}$ 的%
\trCompactness{}得到，这样我们便证明了这些\trFlow{}映射的确属于 $G$。引理的%
证明至此完成，\cref{thm:minimax-critical-value} 的证明也随之完成。
\end{proof}
